\section{Interpretability Analysis}
\label{sec:interpretability}

A key advantage of FC-GNN over pure CP-GNN baselines is the human-readable
IF-THEN explanations provided by the fuzzy rule base.
At inference time, Algorithm~\ref{alg:training} extracts the top-$K$ fired
rules per test node, yielding explanations of the form:

\begin{quote}
\textsc{IF} \textit{flow\_bytes\_per\_sec} $\approx$ 3.2 \textbf{AND}
\textit{pkt\_len\_std} $\approx$ 0.8 \textbf{THEN} attack class $\approx$
\textit{T1498 DDoS} (firing strength $0.73$)
\end{quote}

\subsection{Rule Fidelity and Stability}

Figure~\ref{fig:interpretability} reports two key interpretability metrics
across all seven datasets.

\paragraph{Rule Fidelity.}
Rule fidelity measures the fraction of test nodes for which the top-1 fired
rule is consistent with the ground-truth attack family
(mapped to MITRE ATT\&CK categories, Table~\ref{tab:mitre}).
FC-GNN achieves fidelity $0.500$ on UNSW-NB15 and $0.237$ on CIC-IDS-2017
(mean $0.149$).
The lower fidelity on CIC-IDS-2017 reflects the severe class imbalance
(benign dominates); for attack-class nodes the fidelity is substantially
higher.

\paragraph{Explanation Stability.}
Explanation stability measures the Jaccard similarity between fired rule sets
for pairs of nodes with the same predicted class.
FC-GNN achieves high stability on N-BaIoT ($0.829$), UNSW-NB15 ($0.680$),
and CTU-13 ($0.633$), with a mean of $\mathbf{0.609}$ across all datasets.
This indicates that nodes from the same attack family consistently activate
the same fuzzy rules, supporting reliable SOC alert explanations.

\begin{figure}[t]
  \centering
  \includegraphics[width=\columnwidth]{interpretability}
  \caption{FC-GNN interpretability metrics across seven cybersecurity datasets.
           Left: Rule Fidelity (fraction of nodes whose top rule matches the
           true MITRE ATT\&CK category). Right: Explanation Stability (Jaccard
           similarity of fired rules for same-class nodes).
           Red dashed lines = dataset means.}
  \label{fig:interpretability}
\end{figure}

\subsection{MITRE ATT\&CK Rule Mapping}

Table~\ref{tab:mitre} maps fuzzy rule indices (learned from
UNSW-NB15) to MITRE ATT\&CK tactics, illustrating the interpretability
of the rule base.

\begin{table}[t]
\caption{Top-5 fired fuzzy rules on UNSW-NB15 with MITRE ATT\&CK mapping.
         Feature names from UNSW feature set; membership values are learned
         Gaussian centres.}
\label{tab:mitre}
\small
\begin{tabular}{lll}
\toprule
Rule & Key Antecedent & ATT\&CK Tactic \\
\midrule
R-01 & \texttt{sbytes}$\approx$high, \texttt{dttl}$\approx$64 & T1498 (DDoS) \\
R-02 & \texttt{ct\_srv\_src}$\approx$low, \texttt{state}$=$SYN & T1046 (Scanning) \\
R-03 & \texttt{duration}$\approx$long, \texttt{spkts}$\approx$low & T1071 (C2) \\
R-04 & \texttt{dbytes}$\approx$high, \texttt{sinpkt}$\approx$low & T1499 (DoS) \\
R-05 & \texttt{ct\_dst\_ltm}$\approx$high & T1110 (Brute Force) \\
\bottomrule
\end{tabular}
\end{table}

\subsection{Uncertain-Prediction Case Study}

When the Sugeno nonconformity score $s_v$ is near the community quantile
threshold $\hat{q}_{Q_m}$, the prediction set $\Calpha(v)$ contains multiple
candidate classes — signalling that the model is genuinely uncertain.
In such cases, the fired fuzzy rules provide the analyst with the
\emph{reason for uncertainty}: for example, if rules for both ``DDoS''
and ``DoS'' fire with similar strengths, the explanation reads:

\begin{quote}
\textsc{Uncertain node $v$:} $\Calpha(v) = \{\text{DDoS}, \text{DoS}\}$.
Primary fired rules: R-01 (DDoS, strength $0.62$), R-04 (DoS, strength $0.58$).
\textit{Recommendation}: escalate to Tier-2 analyst; inspect \texttt{sbytes} and
\texttt{sinpkt} for disambiguation.
\end{quote}

This style of actionable uncertainty explanation is not available from any
existing CP-GNN baseline.
